\section{Scaling Law 的范式演变}

前面讲 GPT 和 ChatGPT 的冲击，这一章进入更底层的技术范式。Scaling Law 一开始像单轴规模问题，后来逐渐变成多轴能力增长问题。

访谈中“Scaling Law 的范式演变”是技术主轴之一。张鹏提到，未来不只是基础模型能力持续提升，还包括多模型、多数据融合、智能体和新的计算方式，尤其 RL 方向会继续产生新方法。也就是说，Scaling Law 不再只等同于“更多参数、更多数据、更多算力”的预训练公式。

\begin{figure}[H]
\centering
\includegraphics[width=0.96\textwidth]{scaling-law-evolution.png}
\caption{Scaling Law 范式演变：规模、多模态、后训练、RL 与 Agent。}
\end{figure}

\begin{knowledgebox}{读图：Scaling Law 正在从单轴变多轴}
早期大模型叙事更强调参数、数据和算力；后续则加入多模态数据、后训练、偏好/任务反馈、RL 和 Agent 环境。能力增长不只来自“更大”，也来自“更会用反馈”。
\end{knowledgebox}

\subsection{Agent 为什么成为落地路径}

张鹏认为智能体很重要，因为它解决模型到实际应用之间的落地路径问题。模型本身是能力底座，但客户要的是业务结果。Agent 把模型和工具、流程、环境、用户需求连接起来，是商业化和能力释放的中间层。

\begin{knowledgebox}{术语消化：Scaling Law 相关}
\begin{tabularx}{\textwidth}{p{0.23\textwidth}p{0.32\textwidth}X}
\toprule
术语 & 一句话解释 & 本期中的作用 \\
\midrule
Scaling Law & 模型能力随参数、数据和算力扩展而变化的规律 & 大模型范式的底层叙事。 \\
Post-training & 预训练后用偏好、任务和反馈继续优化模型 & 让模型更可用、更可控。 \\
RL & Reinforcement Learning，强化学习 & 张鹏认为会带来新的计算方式。 \\
Agent & 调用工具并执行任务的智能体 & 模型连接真实应用的落地路径。 \\
多模型融合 & 多种模型和数据源共同构成能力系统 & 基座模型未来方向之一。 \\
\bottomrule
\end{tabularx}
\end{knowledgebox}

\subsection{本章小结}

Scaling Law 的演变说明，大模型竞争不再只是“谁更大”。反馈、后训练、Agent、RL 和多模态数据，正在成为新的能力增长轴。

